\subsection{Теплоемкость}
%%%%%%%%%%%%%%%%%%%%%%%%%%%%%%%%%%%%%%%%%%%%%%%%%%%%
%%%%%%%%%%%%%%%%%%%%%%%%%%%%%%%%%%%%%%%%%%%%%%%%%%%%
\z{В алюминиевый калориметр, в котором была вода массой $m_1 = 200\,\text{г}$ при температуре $t_1 = 10\,\degree\text{C}$, положили медный брусок массой $m_2 = 150\,\text{г}$, температура которого $t_2 = 100\,\degree\text{C}$. Определите температуру воды после установления теплового равновесия. Масса калориметра $m_3 = 100\,\text{г}$.}
%
{$T_{\text{к}} = 15\,\degree\text{C}$}

%%%%%%%%%%%%%%%%%%%%%%%%%%%%%%%%%%%%%%%%%%%%%%%%%%%%
\z{Температура холодной и горячей воды соответственно $t_1 = 12\,\degree\text{C}$ и $t_2 = 70\,\degree\text{C}$. Сколько холодной и горячей воды потребуется, чтобы наполнить ванну водой при температуре $t = 37\,\degree\text{C}$? Масса воды в ванной $M = 150\,\text{кг}$.}
%
{$m_{\text{г}} = 85\,\text{кг},\quad m_{\text{х}} = 65\,\text{кг}$}
	
%%%%%%%%%%%%%%%%%%%%%%%%%%%%%%%%%%%%%%%%%%%%%%%%%%%%
\z{Сколько нужно смешать горячей воды, имеющей температуру $t_1 = 90\,\degree\text{C}$, и холодной, имеющей температуру $t_2 = 10\,\degree\text{C}$, чтобы получить воду массой $m = 100\,\text{кг}$ и температурой $t = 30\,\degree\text{C}$.}
%
{$m_{\text{г}} = 25\,\text{кг},\quad m_{\text{х}} = 75\,\text{кг}$}

%%%%%%%%%%%%%%%%%%%%%%%%%%%%%%%%%%%%%%%%%%%%%%%%%%%%	
\z{В воду массой $m_1 = 200\,\text{г}$ при температуре $t_1 = 20\,\degree\text{C}$ помещают стальную деталь массой $m_2 = 300\,\text{г}$, имеющую температуру $t_2 = 10\,\degree\text{C}$, и медную пластинку массой $m_3 = 400\,\text{г}$ при температуре $t_3 = 25\,\degree\text{C}$. Найдите установившуюся температуру.}
%
{$T_{\text{к}} = 19.5\,\degree\text{C}$}

%%%%%%%%%%%%%%%%%%%%%%%%%%%%%%%%%%%%%%%%%%%%%%%%%%%%
\z {В калориметр с водой перенесли из кипятка металлический шарик, в результате чего температура в калориметре поднялась с $t_1 = 20\,\degree\text{C}$ до $t_2 = 40\,\degree\text{C}$. Какой станет температура в калориметре после переноса из кипятка второго такого же шарика? Третьего? Сколько таких шариков надо перенести, чтобы температура в калориметре стала равной $t = 90\,\degree\text{C}$?}
%
{$52\,\degree\text{C},\quad 60\,\degree\text{C},\quad 21\,\text{штуки}$}

%%%%%%%%%%%%%%%%%%%%%%%%%%%%%%%%%%%%%%%%%%%%%%%%%%%%
\z {При нагревании отопления чугунная батарея массой $m = 20\,\text{кг}$ нагрелась от $t_1 = 10\,\degree\text{C}$ до $t_2 = 50\,\degree\text{C}$, а воздух в комнате объемом $V = 60\,\text{м}^3$ - от $t_1 = 10\,\degree\text{C}$ до $t_3 = 20\,\degree\text{C}$. Для чего потребовалось больше энергии: для нагревания воздуха или батареи? Во сколько раз больше?}
%
{$\frac{Q_{\text{возд}}}{Q_{\text{чуг}}} = 1.67$}
%

%%%%%%%%%%%%%%%%%%%%%%%%%%%%%%%%%%%%%%%%%%%%%%%%%%%%
\z {Нагреватель мощностью $P = 50\,\text{кВт}$ повышает температуру воды, протекающей со скоростью $v = 1\,\text{м/с}$ по трубе диаметром $d = 2\,\text{см}$, с $t_1 = 15\,\degree\text{C}$ до $t_2 = 35\,\degree\text{C}$. Какая часть количества теплоты, выделяемого нагревателем, передается окружающей среде?}%
{$\eta = 0.4725$}
%



